\chapter*{ABSTRAK}
\addcontentsline{toc}{chapter}{ABSTRAK}

\begin{center}

    \textbf{\MakeUppercase{Pengembangan Sistem Pemantauan Sampah Berbasis \protect\textit{Internet of Things} sebagai Dasar Analisis Pola Pembuangan Sampah}} \\
    \vspace{0.4cm}
    oleh \\
    \vspace{0.4cm}
    Jazmy Izzati Alamsyah\\
    18221124\\
\end{center}


\begin{spacing}{1.0}
Pengelolaan sampah di lingkungan kampus menghadapi tantangan berupa kurangnya data terstruktur mengenai pola pembuangan sampah, sehingga menyulitkan pengambilan keputusan berbasis data dalam jadwal pengangkutan dan alokasi sumber daya. Penelitian ini mengembangkan sistem pemantauan sampah berbasis \textit{Internet of Things} (IoT) yang mampu mengotomasi pengumpulan data pembuangan sampah di Gedung Labtek VIII Institut Teknologi Bandung, khususnya pada area depan Laboratorium IoT Lantai 4, sebagai dasar bagi analisis pola pembuangan sampah. Sistem dirancang menggunakan arsitektur IoT 4-\textit{layer} yang terdiri dari \textit{Perception Layer} (mikrokontroler ESP32, sensor jarak VL53L0X, sensor arus INA219, dan modul GSM SIM800L), \textit{Network Layer} berbasis protokol MQTT, \textit{Middleware Layer} menggunakan Node-RED untuk pemrosesan aliran data, serta \textit{Application Layer} dengan basis data Supabase dan \textit{dashboard} Next.js untuk visualisasi data. Selama periode observasi 37 hari, tiga unit perangkat sensor berhasil mengumpulkan 816 \textit{records} data dengan tingkat keberhasilan pengiriman sebesar 91,9\%, dengan 805 \textit{records} dinyatakan valid setelah \textit{filtering} data kotor. Analisis deskriptif yang dilakukan terhadap data tersebut mengidentifikasi empat pola pembuangan sampah, yang terdiri dari tiga pola temporal dan satu pola komposisi, yaitu pola \textit{peak hour} yang konsisten pada pukul 14:00--17:00 (29,6\% dari total \textit{event} penambahan), pola konsentrasi aktivitas pada jam kerja (07:00--19:00) di hari kerja, pola pengaruh hari libur berupa penurunan aktivitas pembuangan, serta pola dominasi sampah anorganik yang menyumbang 13 dari 27 \textit{event} penambahan (48,1\% dari total \textit{event}) dengan kontribusi volume sebesar 51,7\% dari total delta positif. Perangkat bawaan pihak ketiga diadaptasi melalui lima penyesuaian \textit{firmware}, dengan empat di antaranya terbukti efektif dalam meningkatkan reliabilitas dan konsistensi pengukuran. Seluruh 10 kebutuhan fungsional dan 7 dari 8 kebutuhan nonfungsional sistem berhasil dipenuhi. Penelitian ini menunjukkan bahwa sistem pemantauan berbasis \textit{Internet of Things} (IoT) dapat menjadi dasar yang andal untuk menghasilkan wawasan terstruktur dalam mendukung pengelolaan sampah kampus.

Kata kunci: \textit{smart waste bin}, \textit{Internet of Things}, ESP32, VL53L0X, MQTT, pola pembuangan sampah, analisis temporal.
\end{spacing}